\documentclass[11pt]{article}

\usepackage{amsmath,amssymb,graphicx,booktabs,geometry,hyperref,array,longtable,pdflscape}

\geometry{margin=1in}
\newcolumntype{L}[1]{>{\raggedright\arraybackslash}p{#1}}
\hypersetup{
  colorlinks=true,
  linkcolor=blue,
  urlcolor=blue
}

\title{Supplementary Material for ``A Homogeneous Compensation Mechanism for Indirect Evolutionary Rescue in a Roy-Style Predator--Prey Model''}
\author{Draft author list to be finalized}
\date{\today}

\begin{document}

\maketitle

\section*{S1. Numerical classification protocols}

The main manuscript reports consolidated outputs from the numbered analysis scripts rather than new simulations. Classification used the same vocabulary across ODE and PDE outputs:
\begin{itemize}
  \item \textbf{Persistent}: predator density remains positive and late-window change/residual diagnostics indicate a steady or effectively settled state.
  \item \textbf{Extinct}: predator density approaches the extinction threshold and late-window diagnostics do not indicate recovery.
  \item \textbf{Transient}: the trajectory does not meet persistent or extinct steady criteria within the tested horizon, often because late-window change or residual diagnostics remain non-negligible.
\end{itemize}

The tail-window logic was used to avoid treating long transients as equilibria. For PDE tests, classification of spatial means was paired with spatial coefficient-of-variation diagnostics for \(n\), \(w\), and \(q\). A basin label was compared with a matched homogeneous control when heterogeneous perturbation effects were evaluated. The classification logic is diagnostic, not a proof of global basin structure.

\section*{S2. Mathematical appendix: ODE Jacobian and Routh--Hurwitz margins}

Write
\[
F_n=nR,\quad R=r(q)z-\xi-a(q)w,
\]
\[
F_w=wP,\quad P=b(q)nz-(m+s)-\mu w,
\]
\[
F_q=\nu q(1-q)G,\quad G=\Delta r z-\Delta a w,
\]
with \(z=\kappa^{-1}-n-w\), \(\Delta r=r_v-r_u\), \(\Delta a=a_v-a_u\), and \(\Delta b=b_v-b_u\). Since \(z_n=z_w=-1\),
\[
R_n=-r(q),\qquad R_w=-r(q)-a(q),\qquad R_q=G,
\]
\[
P_n=b(q)(z-n),\qquad P_w=-b(q)n-\mu,\qquad P_q=\Delta b\,nz,
\]
\[
G_n=-\Delta r,\qquad G_w=-\Delta r-\Delta a,\qquad G_q=0.
\]
The full reaction Jacobian in variables \((n,w,q)\) is therefore
\[
J=
\begin{pmatrix}
R-nr(q) & n[-r(q)-a(q)] & nG\\
w\,b(q)(z-n) & P+w[-b(q)n-\mu] & w\,\Delta b\,nz\\
-\nu q(1-q)\Delta r & -\nu q(1-q)(\Delta r+\Delta a) & \nu(1-2q)G
\end{pmatrix}.
\]
At an interior compensation-branch equilibrium, \(R=P=G=0\), giving the simplified branch Jacobian
\[
J_*=
\begin{pmatrix}
-n^*r(q^*) & -n^*[r(q^*)+a(q^*)] & 0\\
w^*b(q^*)(z^*-n^*) & -w^*[b(q^*)n^*+\mu] & w^*\Delta b\,n^*z^*\\
-\nu q^*(1-q^*)\Delta r & -\nu q^*(1-q^*)(\Delta r+\Delta a) & 0
\end{pmatrix}.
\]
For \(\det(\lambda I-J)=\lambda^3+A_1\lambda^2+A_2\lambda+A_3\), the coefficients were computed as
\[
A_1=-\mathrm{tr}(J),\qquad
A_2=\frac{1}{2}\left[(\mathrm{tr}J)^2-\mathrm{tr}(J^2)\right],\qquad
A_3=-\det(J).
\]
Table~\ref{tab:rh-margins} reports the Routh--Hurwitz terms and the minimum margin at the target stresses.

\begin{table}[htbp]
  \centering
  \caption{Routh--Hurwitz margins for the linear compensation branch at target stresses. Values are copied to \texttt{results/roy\_ode\_compensation\_routh\_hurwitz\_margins.csv}.}
  \label{tab:rh-margins}
  \small
  \begin{tabular}{rrrrrr}
    \toprule
    \(s\) & \(A_1\) & \(A_2\) & \(A_3\) & \(A_1A_2-A_3\) & min margin \\
    \midrule
    0 & 3.9468 & 0.33673 & 0.003540 & 1.3255 & 0.003540 \\
    0.069448242 & 4.3241 & 0.27528 & 0.004006 & 1.1863 & 0.004006 \\
    0.11765625 & 4.5861 & 0.21486 & 0.003566 & 0.98178 & 0.003566 \\
    0.1584375 & 4.8077 & 0.15238 & 0.002706 & 0.72989 & 0.002706 \\
    0.16486816 & 4.8426 & 0.14158 & 0.002529 & 0.68308 & 0.002529 \\
    0.175 & 4.8976 & 0.12403 & 0.002228 & 0.60525 & 0.002228 \\
    \bottomrule
  \end{tabular}
\end{table}

\begin{figure}[htbp]
  \centering
  \includegraphics[width=0.92\linewidth]{figures/roy_evo_spatial/report/fig42_routh_hurwitz_terms.png}
  \caption{Full Routh--Hurwitz coefficient diagnostics for the linear compensation branch. The main text shows the minimum margin and maximum eigenvalue real part; this supporting figure keeps the individual \(A_1\), \(A_2\), \(A_3\), and \(A_1A_2-A_3\) terms visible for reproducibility.}
  \label{fig:supp-rh-terms}
\end{figure}

\section*{S3. Additional ODE basin maps}

The main text focuses on the derived compensation branch and its stability. Additional ODE diagnostics support the basin interpretation and the homogeneous reaction-level mechanism. Figure~\ref{fig:supp-ode-timeseries} shows representative ODE time series and selection/growth diagnostics. Figure~\ref{fig:supp-ode-basin-map} shows the ODE \(q_0\)--\(w_0\) basin map used to relate initial defense and predator density to persistent, extinct, and transient outcomes. Figure~\ref{fig:supp-ode-phase-equilibria} shows selection-growth phase diagnostics and numerical equilibrium stability.

\begin{figure}[htbp]
  \centering
  \includegraphics[width=0.88\linewidth]{figures/roy_evo_spatial/report/fig28_ode_representative_timeseries.png}
  \caption{Representative ODE trajectories for predator density, defense frequency, selection-gradient diagnostics, and predator-growth diagnostics. These panels support the homogeneous mechanism interpretation but do not replace the analytic branch derivation.}
  \label{fig:supp-ode-timeseries}
\end{figure}

\begin{figure}[htbp]
  \centering
  \includegraphics[width=0.82\linewidth]{figures/roy_evo_spatial/report/fig29_ode_basin_map.png}
  \caption{ODE basin map on the targeted \(q_0\)--\(w_0\) grid. This figure supports basin dependence in the homogeneous reaction system and is not an exhaustive basin-boundary refinement.}
  \label{fig:supp-ode-basin-map}
\end{figure}

\begin{figure}[htbp]
  \centering
  \begin{minipage}{0.49\linewidth}
    \centering
    \includegraphics[width=\linewidth]{figures/roy_evo_spatial/report/fig30_selection_growth_phase.png}
  \end{minipage}
  \hfill
  \begin{minipage}{0.49\linewidth}
    \centering
    \includegraphics[width=\linewidth]{figures/roy_evo_spatial/report/fig31_ode_equilibria_stability.png}
  \end{minipage}
  \caption{Additional homogeneous mechanism diagnostics. Left: selection-gradient and predator-growth relationships for representative cases. Right: numerical equilibria and stability labels across tested stresses.}
  \label{fig:supp-ode-phase-equilibria}
\end{figure}

\section*{S4. PDE non-homogeneous perturbation diagnostics}

The main manuscript reports the conclusion from spatial-mode stability and long-horizon follow-up. The supporting PDE diagnostics are shown here. Figure~\ref{fig:supp-pde-linear} summarizes additional spatial-stability views. Figure~\ref{fig:supp-pde-lambda} shows the continuous \(\lambda\)-scan over the mode-equivalent range \(i,j=0,\ldots,64\), sampled at 600 equally spaced \(\lambda\) values. Figures~\ref{fig:supp-pde-initial-fields} and~\ref{fig:supp-pde-final-fields} show representative non-homogeneous initial and final fields. Figures~\ref{fig:supp-pde-timeseries} and~\ref{fig:supp-pde-cv} show mean dynamics and spatial coefficient-of-variation diagnostics for the targeted non-homogeneous tests. Figures~\ref{fig:supp-pde-long-horizon-timeseries}--\ref{fig:supp-pde-long-horizon-fields} show the long-horizon follow-up diagnostics for the finite-horizon basin-changing cases.

\begin{figure}[htbp]
  \centering
  \begin{minipage}{0.49\linewidth}
    \centering
    \includegraphics[width=\linewidth]{figures/roy_evo_spatial/report/fig46_pde_spatial_stability_along_branch.png}
  \end{minipage}
  \hfill
  \begin{minipage}{0.49\linewidth}
    \centering
    \includegraphics[width=\linewidth]{figures/roy_evo_spatial/report/fig47_pde_diffusion_ratio_stability.png}
  \end{minipage}
  \caption{Additional linear PDE spatial-stability diagnostics. These figures support stability for the tested branch states and diffusion-ratio grid, not a theorem over all diffusion coefficients.}
  \label{fig:supp-pde-linear}
\end{figure}

\begin{figure}[htbp]
  \centering
  \includegraphics[width=0.78\linewidth]{figures/roy_evo_spatial/report/fig66_pde_continuous_lambda_scan.png}
  \caption{Continuous \(\lambda\)-scan for \(J_F(U^*)-\lambda D\) over 600 equally spaced \(\lambda\) values in the mode-equivalent range corresponding to \(i,j=0,\ldots,64\). This scan strengthens the tested spatial-mode result over that range, but it is still a tested-range diagnostic rather than a proof for all possible modes or diffusion coefficients.}
  \label{fig:supp-pde-lambda}
\end{figure}

\begin{figure}[p]
  \centering
  \includegraphics[width=0.96\linewidth]{figures/roy_evo_spatial/report/fig49_nonhomogeneous_initial_conditions.png}
  \caption{Representative non-homogeneous initial PDE fields for targeted finite-amplitude perturbations. This diagnostic figure documents the tested perturbation classes and is not an exhaustive set of heterogeneous initial conditions.}
  \label{fig:supp-pde-initial-fields}
\end{figure}

\begin{figure}[p]
  \centering
  \includegraphics[width=0.96\linewidth]{figures/roy_evo_spatial/report/fig50_nonhomogeneous_final_fields.png}
  \caption{Final PDE fields for representative targeted non-homogeneous perturbation tests. This diagnostic figure supports the absence of persistent spatial patterning in the displayed cases, but it is not a proof for all heterogeneous initial conditions.}
  \label{fig:supp-pde-final-fields}
\end{figure}

\begin{figure}[p]
  \centering
  \includegraphics[width=0.96\linewidth]{figures/roy_evo_spatial/report/fig51_nonhomogeneous_mean_timeseries.png}
  \caption{Mean dynamics for targeted finite-amplitude non-homogeneous PDE tests. These diagnostic trajectories support comparison with matched homogeneous controls; they do not establish global basin behavior.}
  \label{fig:supp-pde-timeseries}
\end{figure}

\begin{figure}[p]
  \centering
  \includegraphics[width=0.96\linewidth]{figures/roy_evo_spatial/report/fig52_nonhomogeneous_spatial_metrics.png}
  \caption{Spatial coefficient-of-variation diagnostics for targeted finite-amplitude non-homogeneous PDE tests. These diagnostics track heterogeneity decay in the tested runs and are not a theorem for all finite-amplitude perturbations.}
  \label{fig:supp-pde-cv}
\end{figure}

\begin{figure}[p]
  \centering
  \includegraphics[width=0.96\linewidth]{figures/roy_evo_spatial/report/fig54_long_horizon_mean_timeseries.png}
  \caption{Long-horizon mean time series for the finite-horizon basin-changing non-homogeneous cases. The followed cases are the targeted cases selected from the earlier finite-amplitude tests, not a broad PDE scan.}
  \label{fig:supp-pde-long-horizon-timeseries}
\end{figure}

\begin{figure}[p]
  \centering
  \includegraphics[width=0.96\linewidth]{figures/roy_evo_spatial/report/fig55_long_horizon_spatial_metrics.png}
  \caption{Long-horizon spatial heterogeneity metrics for the followed basin-changing cases. These diagnostics support spatial-pattern decay in the followed cases.}
  \label{fig:supp-pde-long-horizon-metrics}
\end{figure}

\begin{figure}[p]
  \centering
  \includegraphics[width=0.96\linewidth]{figures/roy_evo_spatial/report/fig56_long_horizon_final_fields.png}
  \caption{Long-horizon final fields for the followed basin-changing cases. These diagnostics support resolution without persistent spatial patterning for the followed cases.}
  \label{fig:supp-pde-long-horizon-fields}
\end{figure}

\section*{S5. Nonlinear trade-off shape-grid details}

The nonlinear extension uses endpoint-preserving shape functions \(f_\gamma(q)=q^\gamma\) with \(\gamma\in\{0.5,1,2\}\). The purpose is to test whether the compensation mechanism survives selected concave, convex, and mixed trade-off shapes, not to scan all possible nonlinear functions. Figure~\ref{fig:supp-nonlinear-branches} shows selected nonlinear branch curves and shape-grid summaries. Figures~\ref{fig:supp-nonlinear-ode-basin}--\ref{fig:supp-nonlinear-pde-heterogeneity} show selected nonlinear ODE basin maps, PDE spatial-stability diagnostics, and non-homogeneous PDE diagnostics.

\begin{figure}[htbp]
  \centering
  \begin{minipage}{0.49\linewidth}
    \centering
    \includegraphics[width=\linewidth]{figures/roy_evo_spatial/report/fig59_nonlinear_branch_curves.png}
  \end{minipage}
  \hfill
  \begin{minipage}{0.49\linewidth}
    \centering
    \includegraphics[width=\linewidth]{figures/roy_evo_spatial/report/fig60_nonlinear_shape_grid_summary.png}
  \end{minipage}
  \caption{Nonlinear branch curves and shape-grid summary. These figures support selected nonlinear compensation branches within the controlled grid.}
  \label{fig:supp-nonlinear-branches}
\end{figure}

\begin{figure}[p]
  \centering
  \includegraphics[width=0.90\linewidth]{figures/roy_evo_spatial/report/fig61_nonlinear_selected_ode_basin_maps.png}
  \caption{Selected nonlinear ODE basin maps. This diagnostic figure supports the controlled nonlinear extension for selected shapes and is not an exhaustive nonlinear basin scan.}
  \label{fig:supp-nonlinear-ode-basin}
\end{figure}

\begin{figure}[p]
  \centering
  \includegraphics[width=0.96\linewidth]{figures/roy_evo_spatial/report/fig62_nonlinear_pde_spatial_stability.png}
  \caption{Selected nonlinear PDE spatial-stability diagnostics. This supporting figure reports targeted spatial-stability checks for selected nonlinear branches and does not prove stability for all nonlinear trade-off shapes.}
  \label{fig:supp-nonlinear-pde-stability}
\end{figure}

\begin{figure}[p]
  \centering
  \includegraphics[width=0.96\linewidth]{figures/roy_evo_spatial/report/fig63_nonlinear_nonhomogeneous_pde_tests.png}
  \caption{Selected nonlinear non-homogeneous PDE diagnostics. This supporting figure reports targeted finite-amplitude checks for selected nonlinear branches and does not establish global robustness to heterogeneous initial conditions.}
  \label{fig:supp-nonlinear-pde-heterogeneity}
\end{figure}

\section*{S6. Additional figure list and CSV source mapping}

\begin{landscape}

Table~\ref{tab:script-source-map} maps the main analysis scripts to the CSV outputs used by the manuscript package. Table~\ref{tab:claim-source-map} maps the main claims to source CSVs and representative figures. The status column records whether each item is a main result or supporting diagnostic evidence. All script names in these tables are relative to \texttt{experiments/}; all CSV names are relative to \texttt{results/}. For readability, the tables use numbered script short names and compact CSV names; the full paths are present in the repository.

\scriptsize
\begin{longtable}{L{0.25\linewidth}L{0.65\linewidth}}
  \caption{Script-to-CSV mapping for reproducibility.}\label{tab:script-source-map}\\
  \toprule
  Script & CSV outputs used here \\
  \midrule
  \endfirsthead
  \toprule
  Script & CSV outputs used here \\
  \midrule
  \endhead
  Step 23: compensation conditions & \texttt{compensation\_conditions\_summary}; \texttt{compensation\_condition\_grid} \\
  Step 24: Routh--Hurwitz & \texttt{routh\_hurwitz\_current}; \texttt{routh\_hurwitz\_summary} \\
  Step 25: PDE spatial and heterogeneous tests & \texttt{pde\_spatial\_modes\_current}; \texttt{pde\_spatial\_nonlinear\_decision} \\
  Step 26: long-horizon follow-up & \texttt{pde\_long\_horizon\_decision}; summary and time-series CSVs \\
  Step 27: nonlinear trade-offs & \texttt{nonlinear\_compensation\_decision}; \texttt{nonlinear\_shape\_summary}; \texttt{nonlinear\_branch\_recovery} \\
  Step 28: counterfactual and modal scan & \texttt{no\_evolution\_counterfactual\_summary}; stress-sweep and time-series CSVs; \texttt{pde\_lambda\_scan\_summary}; \texttt{rh\_margins} \\
  \bottomrule
\end{longtable}

\scriptsize
\begin{longtable}{L{0.28\linewidth}L{0.33\linewidth}L{0.22\linewidth}L{0.10\linewidth}}
  \caption{Claim-to-source mapping for the manuscript package.}\label{tab:claim-source-map}\\
  \toprule
  Claim & Source CSV & Representative figure & Status \\
  \midrule
  \endfirsthead
  \toprule
  Claim & Source CSV & Representative figure & Status \\
  \midrule
  \endhead
  Frozen-\(q\) counterfactual verifies indirect rescue over a finite stress window & \texttt{no\_evolution\_counterfactual\_summary}; stress-sweep CSV & Main no-evolution figure & Main claim \\
  Linear compensation branch conditions are satisfied in the current parameterization & \texttt{compensation\_conditions\_summary} & Main branch figure; Fig.~\ref{fig:supp-ode-phase-equilibria} & Main claim \\
  Routh--Hurwitz stability holds at target stresses & \texttt{routh\_hurwitz\_summary} & Main Routh--Hurwitz figure; Fig.~\ref{fig:supp-rh-terms} & Main claim \\
  Spatial modes are stable for the current PDE coefficients over the tested range & \texttt{pde\_spatial\_nonlinear\_decision}; \texttt{pde\_lambda\_scan\_summary} & Main spatial-stability figure; Fig.~\ref{fig:supp-pde-linear}; Fig.~\ref{fig:supp-pde-lambda} & Main claim \\
  Long-horizon non-homogeneous basin changes resolve to homogeneous controls & \texttt{pde\_long\_horizon\_decision} & Main long-horizon decision panel; Fig.~\ref{fig:supp-pde-long-horizon-timeseries}; Fig.~\ref{fig:supp-pde-long-horizon-metrics}; Fig.~\ref{fig:supp-pde-long-horizon-fields} & Main claim \\
  Linear branch is recovered by nonlinear branch formulation & \texttt{nonlinear\_branch\_recovery} & Fig.~\ref{fig:supp-nonlinear-branches} & Supporting \\
  Nonlinear shape grid includes robust, partial, no-compensation, and unresolved cases & \texttt{nonlinear\_shape\_summary}; \texttt{nonlinear\_compensation\_decision} & Main nonlinear heatmap; Fig.~\ref{fig:supp-nonlinear-branches}; Fig.~\ref{fig:supp-nonlinear-ode-basin} & Main claim \\
  \bottomrule
\end{longtable}
\normalsize

\end{landscape}

\begin{figure}[htbp]
  \centering
  \fbox{\parbox{0.86\linewidth}{\centering The main manuscript contains the primary compensation branch and stress-interval figure. This boxed note avoids duplicating that main-text figure in the supplement while keeping the claim-to-source map readable.}}
  \caption{Pointer to the main compensation branch figure.}
  \label{fig:supp-main-branch-note}
\end{figure}

\end{document}
